\documentclass[preview]{standalone}

\usepackage{amsmath}
\usepackage{amssymb}
\usepackage{stellar}
\usepackage{definitions}

\begin{document}

\id{psicologia-condizionamento-classico}
\genpage

\section{Condizionamento classico}

\begin{snippetexample}{film-horror-condizionamento-classico-example}{Film horror}
    Guardando un film horror per la seconda volta, il cuore può iniziare a
    battere più forte appena la colonna sonora suggerisce che sta per accadere
    qualcosa di pericoloso. Il corpo ha imparato a produrre una risposta
    fisiologica quando un evento situazionale, come una musica carica di
    tensione, è associato a uno stimolo visivo spaventoso.
\end{snippetexample}

\begin{snippetdefinition}{condizionamento-classico-definition}{Condizionamento classico}
    Il \emph{condizionamento classico} è una forma di apprendimento in cui uno
    stimolo o evento predice il verificarsi di un secondo stimolo o evento,
    permettendo all'organismo di apprendere una nuova associazione tra due
    stimoli.
\end{snippetdefinition}

\begin{snippetexample}{pavlov-condizionamento-classico-example}{Esperimento di Pavlov}
    Ivan Pavlov studiava la digestione nei cani inserendo fistole nelle
    ghiandole salivari e negli organi digestivi, così da raccogliere e misurare
    le secrezioni. Dopo ripetute presentazioni del cibo, osservò che i cani
    iniziavano a salivare prima di essere nutriti, per esempio alla vista
    dell'assistente o al rumore dei suoi passi. Negli esperimenti successivi,
    uno stimolo acustico inizialmente neutro precedeva regolarmente la
    presentazione del cibo; dopo varie associazioni, il solo stimolo acustico
    era sufficiente a produrre salivazione.
\end{snippetexample}

\includesnpt[width=72\%,src=/snippet/static-psychology/pavlov-classical-conditioning-experiment.jpg]{centered-img}

\begin{snippetdefinition}{stimolo-incondizionato-definition}{Stimolo incondizionato}
    Lo \emph{stimolo incondizionato}, o \emph{SI}, è uno stimolo che attiva
    naturalmente un riflesso fisiologico, senza necessità di apprendimento.
\end{snippetdefinition}

\begin{snippetdefinition}{risposta-incondizionata-definition}{Risposta incondizionata}
    La \emph{risposta incondizionata}, o \emph{RI}, è il comportamento attivato
    naturalmente da uno stimolo incondizionato.
\end{snippetdefinition}

\begin{snippetdefinition}{stimolo-condizionato-definition}{Stimolo condizionato}
    Lo \emph{stimolo condizionato}, o \emph{SC}, è uno stimolo inizialmente
    neutro che diventa capace di attivare un comportamento dopo essere stato
    associato regolarmente a uno stimolo incondizionato.
\end{snippetdefinition}

\begin{snippetdefinition}{risposta-condizionata-definition}{Risposta condizionata}
    La \emph{risposta condizionata}, o \emph{RC}, è un comportamento attivato da
    uno stimolo inizialmente neutro diventato rilevante attraverso
    l'associazione con uno stimolo incondizionato.
\end{snippetdefinition}

\begin{snippet}{connessioni-condizionamento-classico}
    Nel condizionamento classico, la biologia fornisce le connessioni
    \texttt{SI-RI}, mentre l'apprendimento crea le connessioni
    \texttt{SC-RC}. Prima del condizionamento, lo stimolo neutro non produce la
    risposta; durante il condizionamento viene abbinato allo stimolo
    incondizionato; dopo il condizionamento produce una risposta condizionata.
\end{snippet}

\begin{snippetdefinition}{acquisizione-condizionamento-definition}{Acquisizione}
    L'\emph{acquisizione} è la fase del condizionamento classico in cui si forma
    l'associazione tra stimolo condizionato e stimolo incondizionato, con un
    aumento progressivo della risposta condizionata.
\end{snippetdefinition}

\begin{snippetdefinition}{estinzione-condizionamento-definition}{Estinzione}
    L'\emph{estinzione} è la diminuzione della risposta condizionata che avviene
    quando lo stimolo condizionato viene presentato ripetutamente senza lo
    stimolo incondizionato.
\end{snippetdefinition}

\begin{snippetdefinition}{recupero-spontaneo-definition}{Recupero spontaneo}
    Il \emph{recupero spontaneo} è la ricomparsa di una risposta condizionata
    dopo un periodo di riposo, anche in assenza dello stimolo incondizionato.
\end{snippetdefinition}

\begin{snippetdefinition}{riacquisizione-condizionamento-definition}{Riacquisizione}
    La \emph{riacquisizione} è il rapido ritorno della risposta condizionata
    quando lo stimolo condizionato viene nuovamente presentato insieme allo
    stimolo incondizionato dopo una fase di estinzione.
\end{snippetdefinition}

\includesnpt[width=58\%,src=/snippet/static-psychology/conditioning-acquisition-extinction-reacquisition.jpg]{centered-img}

\begin{snippet}{timing-condizionamento-classico}
    Nel condizionamento classico la contiguità temporale tra stimolo
    condizionato e stimolo incondizionato è importante. Le principali modalità
    di presentazione sono:
    \begin{itemize}
        \item \emph{condizionamento ritardato}: lo stimolo condizionato compare
        prima dello stimolo incondizionato e rimane fino alla sua comparsa;
        \item \emph{condizionamento di traccia}: lo stimolo condizionato
        scompare prima dello stimolo incondizionato;
        \item \emph{condizionamento simultaneo}: i due stimoli compaiono nello
        stesso momento;
        \item \emph{condizionamento retrogrado}: lo stimolo condizionato compare
        dopo lo stimolo incondizionato.
    \end{itemize}
\end{snippet}

\includesnpt[width=78\%,src=/snippet/static-psychology/conditioning-timing-procedures.jpg]{centered-img}

\begin{snippetdefinition}{generalizzazione-stimolo-definition}{Generalizzazione dello stimolo}
    La \emph{generalizzazione dello stimolo} è l'estensione della risposta
    condizionata a stimoli mai associati allo stimolo incondizionato originario,
    ma simili allo stimolo condizionato iniziale.
\end{snippetdefinition}

\includesnpt[width=45\%,src=/snippet/static-psychology/stimulus-generalization-curves.jpg]{centered-img}

\begin{snippetdefinition}{discriminazione-stimolo-definition}{Discriminazione dello stimolo}
    La \emph{discriminazione dello stimolo} è il processo attraverso cui un
    organismo impara a rispondere in modo diverso a stimoli distinti dallo
    stimolo condizionato per qualche dimensione.
\end{snippetdefinition}

\begin{snippet}{generalizzazione-discriminazione}
    Generalizzazione e discriminazione devono trovare un equilibrio. La
    generalizzazione estende l'apprendimento a stimoli simili a quelli già
    incontrati; la discriminazione limita la risposta ai soli stimoli rilevanti,
    evitando reazioni inutili o costose.
\end{snippet}

\begin{snippetexample}{rescorla-contingenza-example}{Esperimento di Rescorla}
    Robert Rescorla mostrò che il condizionamento classico non dipende soltanto
    dalla contiguità temporale. In un esperimento con cani, usò un suono come
    stimolo condizionato e una scossa elettrica come stimolo incondizionato. Un
    gruppo riceveva suono e scossa solo in modo temporalmente vicino; per
    l'altro gruppo, invece, il suono prediceva in modo affidabile la scossa. La
    paura condizionata risultava più forte quando lo stimolo condizionato era
    informativo rispetto allo stimolo incondizionato.
\end{snippetexample}

\begin{snippetexample}{piccolo-albert-example}{Piccolo Albert}
    Watson e Rayner addestrarono Albert ad avere paura di un ratto bianco,
    inizialmente gradito, associando la sua comparsa a un forte rumore
    spiacevole. Il rumore produceva naturalmente paura e stress emotivo; dopo
    l'associazione ripetuta, il ratto bianco divenne capace di suscitare una
    reazione di paura. L'esperimento voleva mostrare che alcune risposte emotive
    possono essere apprese tramite condizionamento.
\end{snippetexample}

\begin{snippetdefinition}{salienza-stimolo-definition}{Salienza dello stimolo}
    La \emph{salienza dello stimolo} è il grado con cui uno stimolo viene
    notato rapidamente, in funzione della sua intensità e del contrasto con gli
    altri stimoli presenti nell'ambiente.
\end{snippetdefinition}

\end{document}
